% Avsnitt 5.1-5.5, ur lectures/L05/appendix/a_variables_and_hardware.md (A.1-A.4).

\section{Vad som återstår}
\label{c5:sec:intro}

Varje konstruktion du behöver för att beskriva hårdvara har du redan använt: \code{entity} och
\code{architecture} i \chapref{c1:ch}, \code{std_logic_vector}, \code{process} och \code{case} i
\chapref{c2:ch}, den klockade processen och schemaläggningsregeln i \chapref{c3:ch},
subkomponenter och generics i \chapref{c2:ch} och \ref{c4:ch}.

Två saker återstår, och de är de två svåraste att lista ut på egen hand:
\begin{itemize}
  \item Den första är \code{variable}, den enda konstruktion boken inte har behövt förrän nu. Den
    ser ut som en signal och beter sig inte alls som en, och skillnaden är precis den sort som ger
    en design som simulerar trovärdigt och är tyst felaktig. \Secref{c5:sec:variable} går igenom
    fallet där det biter.
  \item Den andra är vad något av det här faktiskt blir. Sedan \chapref{c1:ch} har du litat på att
    syntesverktyget bygger den krets du beskrev. Det gör det, men inte av något du skulle känna
    igen. \Secref{c5:sec:fabric} öppnar lådan. Det visar sig att en FPGA inte innehåller några
    grindar alls, vilket är skälet till att \chapref{c2:ch}:s Karnaughdiagram köpte dig mindre än
    de såg ut att göra. Din klocka har en hastighetsgräns (\secref{c5:sec:fmax}), och något
    specifikt sätter den. Och ett \code{case} med en gren som saknas ger dig tyst ett minneselement
    som ingen bad om (\secref{c5:sec:latch}).
\end{itemize}

Inget av avsnitten svarar på ”är min logik rätt?” Den frågan har du besvarat sedan \chapref{c1:ch}.
De svarar på de två frågorna efter den: är det här vad jag menade, och kommer det att få plats och
gå.

% ------------------------------------------------------------------------------------------
\section{\code{variable}: processlokal lagring som uppdateras direkt}
\label{c5:sec:variable}

En \code{variable} deklareras inuti en \code{process}, mellan \code{process(...)} och \code{begin}.
Den avgörande skillnaden mot en \code{signal}: variabeltilldelning använder \code{:=} och får
verkan \textbf{omedelbart}, så redan nästa sats som läser den ser det uppdaterade värdet, precis
som en tilldelning i C. Det finns ingen schemaläggning, ingenting av den
vänta-tills-processen-suspenderar-fördröjning som \secref{c3:sec:template} stavade ut för signaler,
och inget ”sista tilldelningen vinner”.

\begin{booktable}{@{}p{9em}p{9.6em}L@{}}
   & \thead{\code{signal} (\code{<=})} & \thead{\code{variable} (\code{:=})} \\
  \midrule
  Deklareras & I arkitekturens deklarativa del & I processens deklarativa del \\
  Synlig från & Var som helst i arkitekturen & Endast inuti sin egen process \\
  Uppdateringstillfälle & Schemalagd: tillämpas när processen nästa gång suspenderar & Omedelbar:
    tillämpas före nästa sats \\
  Upprepad tilldelning i ett varv & Bara den sista överlever; tidigare kastas & Var och en får
    verkan, i tur och ordning \\
\end{booktable}

\subsection*{Varför det här faktiskt spelar roll: en XOR-paritetsgenerator}
En entitet som beräknar XOR-pariteten hos en 8-bitars ingång, \code{'1'} om \code{bits} har ett
udda antal ettställda bitar:

\begin{vhdlcode}
entity parity_gen is
    port(bits  : in  std_logic_vector(7 downto 0);
         parity: out std_logic);
end entity;
\end{vhdlcode}

Med bakgrund i mjukvara är den första ingivelsen att ackumulera den löpande XOR-summan i en loop,
med en signal som ackumulator:

\begin{vhdlcode}
-- Don't do this: shown to demonstrate why it doesn't work.
architecture broken of parity_gen is
signal parity_s: std_logic;
begin
    parity <= parity_s;

    process(bits) is
    begin
        parity_s <= '0';
        for i in bits'range loop
            parity_s <= parity_s xor bits(i);
        end loop;
    end process;
end architecture;
\end{vhdlcode}

Spåra ett varv, med schemaläggningsregeln från \secref{c3:sec:template}: varje läsning av
\code{parity_s} returnerar det den höll \emph{innan det här varvet började} (kalla det
\code{old_parity}), och bara den sista tilldelningen överlever.
\begin{itemize}
  \item \code{parity_s <= '0';} schemalägger \code{'0'}. Ännu inte tillämpad.
  \item \code{i = 7} (\code{bits'range} på \code{(7 downto 0)} går från hög till låg): läser
    \code{parity_s}, får \code{old_parity}, eftersom \code{'0'}:an inte har tillämpats.
    Schemalägger \code{old_parity xor bits(7)}, och kastar \code{'0'}:an.
  \item \code{i = 6}: läser \code{parity_s}, återigen \code{old_parity}. Schemalägger
    \code{old_parity xor bits(6)}, och kastar den föregående.
  \item \dots och så vidare för varje iteration \dots
  \item \code{i = 0}, den sista satsen i varvet: läser fortfarande \code{old_parity}, schemalägger
    \code{old_parity xor bits(0)}, och kastar allt före den.
  \item Processen suspenderar och \code{parity_s} blir \code{old_parity xor bits(0)}, den enda av
    nio schemalagda tilldelningar som överlevde.
\end{itemize}

Resultatet beror bara på \code{bits(0)} och på vad \code{parity_s} råkade hålla sedan tidigare. Sju
iterationer och den inledande \code{'0'}:an beräknades och slängdes. Det här är inget randfall: det
är vad en signal-som-ackumulator \emph{alltid} gör, eftersom varje iteration läser samma frusna
värde och bara den textuellt sista tilldelningen behålls.

Samma beräkning med en \code{variable}, och det här är det genomarbetade exemplet på disk:

\vhdlfile{lectures/L05/parity_gen/parity_gen.vhd}

Nu tillämpas \code{acc := '0'} omedelbart; \code{i = 7} läser \code{'0'}, XOR:ar in \code{bits(7)}
och uppdaterar \code{acc} på en gång; \code{i = 6} läser det just uppdaterade värdet, och så vidare
genom varje iteration. Loopen slutar med att \code{acc} håller den sanna XOR:en av alla åtta
bitarna, och \code{parity_s <= acc;} schemalägger det enda korrekta värdet, den enda tilldelningen
till \code{parity_s} i varvet.

\subsection*{Vad den loopen inte är}
Det är inte åtta steg som hårdvaran utför ett efter ett. Det är här en bakgrund i C vilseleder som
mest, så det är värt att säga rakt ut: en \code{for}-loop vars gränser är fastlagda vid
elaboreringen \textbf{rullas ut} av syntesverktyget. De åtta iterationerna blir åtta XOR-grindar i
en kedja, som alla lägger sig inom en fortplantningsfördröjning, i samma klockcykel. Det finns
ingen räknare för \code{i}, ingen förgrening, ingen iteration i tiden. Loopen skriver åtta grindar
på tre rader; den gör inte åtta saker i följd.

Det är också därför \code{acc}, som läser sig som en löpande summa, inte kostar någon lagring alls.
Den behöver aldrig överleva en klockflank: vart och ett av dess värden är bara ledningen mellan två
av de där XOR-grindarna. En \code{variable} är rätt verktyg just för att den aldrig överlever ett
varv.

Detsamma gäller \code{bits'range}. \textbf{Attributet} \code{'range} ger indexintervallet för
\code{bits}, här \code{7 downto 0}, så att loopen täcker varje bit utan att gränserna skrivs två
gånger. Det löses upp vid elaboreringen, som allt annat i loopens huvud, vilket är precis därför
verktyget kan rulla ut den. \code{for i in 0 to 1 loop} i \chapref{c3:ch}:s \file{led_toggle.vhd}
betyder samma sak: båda bitarna får sina egna grindar och båda uppdateras i samma cykel.

\subsection*{Tumregeln}
\begin{itemize}
  \item Använd en \code{variable} för en tillfällig beräkning som konsumeras helt inom ett varv:
    loopackumulatorer, temporära byten, mellanvärden som lämnas över till en signal på slutet,
    precis som \code{acc}.
  \item Använd en \code{signal} för allt som en annan process eller omvärlden läser, och allt som
    måste hålla sitt värde över skilda varv (en räknare, eller ett skiftregisters lagrade bitar,
    båda i \chapref{c6:ch}).
\end{itemize}

En variabel \emph{kan} behålla sitt värde mellan varv genom samma process, och beter sig då som
statisk lokal lagring snarare än en ny lokal variabel varje gång, så det går att bygga tillstånd
med en. Men eftersom dess uppdateringar är omedelbara snarare än schemalagda beter den sig
annorlunda än en signal som används på samma sätt, särskilt för allt som modellerar en kedja av
register som stegar ett steg per klockflank. Om tillstånd genuint behöver bestå och läsas av annan
logik, grip efter en \code{signal}: semantiken med schemalagd uppdatering är det som får
\chapref{c6:ch}:s räknare och skiftregister att fungera.

% ------------------------------------------------------------------------------------------
\section{Vad FPGA:n faktiskt består av}
\label{c5:sec:fabric}

En FPGA är inget oskrivet blad som grindar etsas in i. Den är ett fast rutnät av identiska små
block, tillverkade innan någon skrev din VHDL, och ”syntes” är att lista ut hur de ska konfigureras
så att de beter sig som den krets du beskrev. Två av dem spelar roll här:
\begin{itemize}
  \item En \textbf{uppslagstabell}, eller \textbf{LUT}: ett litet minne, typiskt fyra till sex
    ingångar och en utgång. Ladda den med sanningstabellen för vilken boolesk funktion som helst av
    de ingångarna, så beräknar den den funktionen. Det är det ärliga svaret på ”vad blir en grind
    på en FPGA”: ingenting blir en grind. En LUT med fyra ingångar beräknar \code{a and b},
    \code{a xor b or (c and not d)}, eller vilken annan funktion som helst av upp till fyra
    ingångar till exakt samma kostnad, eftersom det i varje fall är samma block som håller en annan
    tabell.
  \item En \textbf{vippa}: en bit klockad lagring, som sitter bredvid varje LUT.
\end{itemize}

På DE0-CV:ns Cyclone~V är de paketerade tillsammans som en \textbf{ALM}, som håller en delbar LUT
med sex ingångar, användbar som en funktion av sex variabler eller två mindre, plus fyra vippor.
Det är den enhet Quartus fitter-rapport räknar, vilket är skälet till att den talar om ALM:ar
snarare än grindar.

Det är nästan hela historien. \code{register4} blir fyra vippor, \code{mux_8to1} en handfull
LUT:ar, och \chapref{c8:ch}:s tillståndsmaskin LUT:ar som matar de vippor som håller tillståndet.

Det förklarar också något från \chapref{c2:ch} som annars skulle se ut som ett brutet löfte: boken
lät dig minimera uttryck med Karnaughdiagram för att spara grindar, och nämnde sedan aldrig
besparingen igen. På en FPGA kostar en funktion av fyra variabler en LUT vare sig du minimerade den
eller inte. Minimering spelar fortfarande roll: det är så du förstår en krets, den är precis rätt
för de flöden med diskret logik och ASIC som Karnaughdiagram uppfanns för, och att reducera en
funktion under LUT:ens ingångsbredd är det som hindrar den från att behöva en andra nivå av LUT:ar.
Men på det här kortet köper den oftast ingenting, och det är bättre att veta det än att undra.

% ------------------------------------------------------------------------------------------
\section{Varför en klocka har en hastighetsgräns}
\label{c5:sec:fmax}

Signaler tar tid på sig genom en LUT och längs ledningarna mellan block. Den
\textbf{fortplantningsfördröjningen} är skälet till att en klocka inte kan gå hur snabbt som helst.

Tänk dig den form varje design i den här boken har, en vippa, lite kombinatorik, ännu en vippa:
\begin{itemize}
  \item På en stigande flank presenterar den första ett nytt värde, som måste färdas genom varje
    LUT och ledning mellan de två och \emph{lägga sig} innan nästa flank anländer.
  \item Om det inte har gjort det fångar den andra vippan ett värde som fortfarande ändras, och
    designen gör fel sak tillförlitligt och gåtfullt.
\end{itemize}

Så klockperioden måste täcka tre saker:
\begin{itemize}
  \item Den första vippans \textbf{clock-to-Q}-fördröjning: tiden mellan flanken och att dess
    utgång faktiskt ändras.
  \item Den \textbf{kombinatoriska fördröjningen} genom LUT:arna och ledningarna mellan de två.
  \item Den andra vippans \textbf{setup-tid}.
\end{itemize}

Den långsammaste sådana vägen är \textbf{kritiska vägen}, och ett genom den summan är \textbf{Fmax}.

Quartus beräknar Fmax åt dig, i tidsanalyssteget (bilaga~\ref{app:quartus}), men bara om du talar
om för den vad klockan är. Det innebär en \file{.sdc}-villkorsfil med ett \code{create_clock} på
\code{50 MHz}-ingången. Utan en sådan har analysatorn ingen period att jämföra något mot, så den
rapporterar varje väg som obegränsad och ger dig ingen Fmax alls. Skilt från det är en design som
”missar timingen” en där någon väg visade sig långsammare än den period du faktiskt deklarerade.

\textbf{Inget av det ingår i den här boken.} Du kommer inte att skriva någon \file{.sdc}, läsa
någon timingrapport, eller möta en design som missar timingen, eftersom varje design här ryms inom
20~ns med marginal. Det är värt att veta att steget finns och vad det kräver, så att du känner igen
det första gången en design av dig närmar sig gränsen.

Två följder värda att bära med sig:
\begin{itemize}
  \item \textbf{Djup kombinatorik kostar dig hastighet, inte antalet vippor.} Att kedja ett långt
    uttryck mellan två register förlänger kritiska vägen; att dela upp det så att en del av arbetet
    sker på en cykel och resten på nästa förkortar den. Den avvägningen, fler register mot en
    snabbare klocka, är det enskilt vanligaste greppet i snabb digital konstruktion.
  \item \textbf{\code{50 MHz} är en budget på 20~ns.} Det är den siffra varje väg i din design
    tävlar mot, och det är gott om tid vid de här hastigheterna.
\end{itemize}

% ------------------------------------------------------------------------------------------
\section{Låset du får av misstag}
\label{c5:sec:latch}

\Chapref{c3:ch} byggde ett D-lås, visade att det var transparent medan det var aktiverat, och lade
det sedan åt sidan till förmån för vippan. Här kommer det tillbaka objudet.

\Secref{c2:sec:process} krävde ett \code{when others} på varje \code{case}, och gav språkets skäl:
ett \code{case} måste täcka varje värde av sin selektors typ. Den allmänna regeln bakom det är att
en kombinatorisk \code{process} måste tilldela sin utgång på \emph{varje} väg genom sig, och det
skälet går nu att formulera:
\begin{itemize}
  \item VHDL kräver att en signal behåller sitt gamla värde om ingenting tilldelar ett nytt.
  \item Om någon väg genom din process lämnar \code{x} otilldelad har du sagt att \code{x} måste
    \emph{minnas} sitt föregående värde på den vägen, och att minnas är inte något ett
    kombinatoriskt nät kan göra.
  \item Så verktyget bygger det enda som kan: ett lås, precis det från \secref{c3:sec:latch}.
\end{itemize}

Det är ett \textbf{oavsiktligt lås}, och det är nästan aldrig vad någon menade. Det är en varning
snarare än ett fel, designen kommer ofta att verka fungera, och meddelandet är värt att känna igen:

\begin{console}
Warning: Inferred latch(es) for signal "x"
\end{console}

Åtgärden är aldrig att lägga till ett lås med flit. Den är att tilldela signalen på varje väg,
vilket är regeln du redan följde.
